%
% listen_lifecycle.tex
%
% A Z specification of luxd's WebSocket listen-leg lifecycle: successive
% sessions of one identity contending for a single connection's
% listener slot, and the menu callbacks whose registration is gated on
% that slot.  Formalises the design that closes the stale-teardown
% clobber and the registration time-of-check/time-of-use gap.
%
% Source of record:
%   src/punt_lux/ws_listen.py                  -- HubListenSession: run,
%                                                 _pump, _teardown
%   src/punt_lux/domain/hub/callback_hold.py   -- CallbackRouter:
%                                                 add_listener,
%                                                 remove_listener,
%                                                 has_listener, route, take
%   src/punt_lux/domain/hub/hub_clients.py     -- HubClientRegistry:
%                                                 record, register_callback,
%                                                 withdraw_callbacks,
%                                                 live_sessions
%   src/punt_lux/operations/callbacks.py       -- CallbackOperations:
%                                                 register_callback (the
%                                                 has_listener gate)
%
% Type-check:  fuzz -t docs/listen_lifecycle.tex
% Model-check (see Section "Verification" for the exact goal predicates):
%   probcli docs/listen_lifecycle.tex -model_check -p DEFAULT_SETSIZE 2
%
% Formatting follows the fuzz house rules: \quad~ for continuation
% indent (never \t1..\t9), schema lines under 80 columns, predicates
% broken at \land/\lor/\implies with the operator starting the
% continuation line.
%

\documentclass[a4paper,10pt,fleqn]{article}
\usepackage[margin=1in]{geometry}
\usepackage{fuzz}
\usepackage[colorlinks=true,linkcolor=blue,citecolor=blue,urlcolor=blue]{hyperref}

\begin{document}

\title{The luxd Listen-Leg Lifecycle: a Z Specification}
\author{Formal model of listener ownership and callback registration}
\date{July 2026}
\maketitle

% ============================================================
\section{Introduction}
% ============================================================

A client holds one persistent WebSocket to \texttt{luxd}.  Over that leg
it receives the pub-sub events it subscribed to and the menu-callback
clicks routed to its session.  The connection is keyed by an
\emph{identity-derived} connection id, so two successive sessions of the
same identity --- an old one dying and a new one reconnecting after a
100ms backoff --- address one and the same slot in the Hub's registries.

That sharing is the whole difficulty.  Each session installs itself as
the connection's listener and writer on connect, and removes them again
on teardown; but a session that has already lost its socket may still be
\emph{suspended} inside its \texttt{finally} block, awaiting a cancelled
writer task, while its successor completes an entire connect.  When the
old session finally resumes, an unconditional teardown removes the
\emph{successor's} state.  The successor keeps its lease alive with
keepalives, so the damage does not heal: a session that is live, that
believes it is push-reachable, and that owns no menu entries.

A second, independent gap lies in registration.  A menu callback may be
registered only by a connection that holds a listen leg, because a menu
item that cannot be pushed to cannot launch at the speed a menu implies.
That gate reads the router under the router's lock and then writes the
registry under the registry's lock.  The two are not one critical
section, so a teardown may pass between them and leave a callback
registered with no listener --- and, since the withdrawal has already
happened, with nothing that will ever withdraw it.

This document models the leg as a finite state machine over a small
fixed set of sessions and callbacks, so that ProB can enumerate every
interleaving and either confirm the three safety invariants or produce
the exact offending trace.

\paragraph{Concurrency model.}
Every session action runs on \texttt{luxd}'s single asyncio loop, so a
run of code between two \texttt{await}s is atomic with respect to other
session actions.  \texttt{run}'s prologue (\texttt{record},
\texttt{register\_writer}, \texttt{add\_listener}) is one such run, and
so is \texttt{\_teardown}, which contains no \texttt{await} at all.
Registration and click routing arrive from \emph{other} threads --- an
MCP tool thread, a REST worker, the click-dispatch thread --- and are
therefore free to interleave anywhere, including between two of
\texttt{\_teardown}'s writes.  The model encodes this asymmetry with a
single component, \texttt{looprun}, which is non-empty exactly while a
loop-side run is in progress and which guards the loop-side operations
and only those.

% ============================================================
\section{Basic Types}
% ============================================================

\begin{zed}
  [SESSION, CALLBACK]
\end{zed}

\texttt{SESSION} is the set of \texttt{HubListenSession} instances that
may occupy the one connection over time.  A session object is
constructed afresh by each run of the \texttt{/ws} route and installs
itself exactly once, so its identity is a monotone stamp: no live
incarnation ever reuses the identity of an earlier one.  That property
is what lets ownership be decided by comparing a stored listener to
\texttt{self}, with no generation counter.  Two sessions suffice to
exhibit both defects; the carrier is bounded by ProB's
\texttt{DEFAULT\_SETSIZE}.

\texttt{CALLBACK} is the set of menu callbacks the connection's app may
register.  Two are enough to let a stale registration and a fresh one
coexist.

% ============================================================
\section{Free Types}
% ============================================================

\begin{zed}
  SPHASE ::= \\
  \quad~ sfresh | sattached | spumping | ssuspended | stearing | sdone
\end{zed}

One session's own progress.  \texttt{sfresh}: constructed, nothing
installed.  \texttt{sattached}: the prologue has run --- the registry
entry is recorded and the writer and listener are installed --- and the
session is suspended on the handshake write.  \texttt{spumping}: the
handshake landed and the read loop is running.  \texttt{ssuspended}: the
read loop has ended (peer gone, protocol error, or a failed handshake
write) and the coroutine is suspended in \texttt{\_pump}'s
\texttt{finally}, awaiting the cancelled writer task --- the window this
model exists to reason about.  \texttt{stearing}: the teardown's
detaching step has fired and the disconnect cascade is still to run.
\texttt{sdone}: finished.

\begin{zed}
  FLAG ::= leaseon | leaseoff
\end{zed}

Whether the connection's registry entry is still within its lease.  Any
recorded contact renews it; it lapses only when no session is attached,
and the sweep that follows takes the entry's callbacks and its hold with
it.

% ============================================================
\section{State}
% ============================================================

The state is flat: one schema, ten components.  Its predicate carries
only \emph{structural} constraints --- cardinality bounds on the
single-occupancy slots, and the disjointness of pending and committed
registrations.  These hold in every design, correct or buggy, so they
act as coherence constraints, not as the safety properties under test.
The safety properties (Section~\ref{sec:inv}) are deliberately
\emph{absent} here, so that a violating state remains reachable and ProB
can find it.

\begin{schema}{Leg}
  phase : SESSION \fun SPHASE \\
  listener : \power SESSION \\
  writer : \power SESSION \\
  incumbent : \power SESSION \\
  owner : CALLBACK \pfun SESSION \\
  gated : CALLBACK \pfun SESSION \\
  held : \power CALLBACK \\
  leased : FLAG \\
  looprun : \power SESSION \\
  clobbered : \power SESSION
\where
  \# listener \leq 1 \\
  \# writer \leq 1 \\
  \# incumbent \leq 1 \\
  \# looprun \leq 1 \\
  \dom gated \cap \dom owner = \emptyset
\end{schema}

\texttt{listener} is the router's registered
\texttt{CallbackListener} for this connection, modelled as a set of size
at most one so that the empty set stands for ``none'' without a
sentinel; \texttt{writer} is the Hub's registered pub-sub writer, held
the same way.  \texttt{owner} maps each \emph{registered} callback to
the session that installed it, so $\dom owner$ is the registry
entry's callback set and $\ran owner$ its provenance;
\texttt{gated} is the same for registrations that have passed the
push-reachability gate but have not yet been written.  \texttt{held} is
the router's hold --- the clicks buffered for this connection.

Two components are ghosts, present only to state properties of the
transition relation as properties of a state.  \texttt{incumbent} names
the session whose connect most recently completed and which has not
itself detached: the session that a correct design keeps
push-reachable.  \texttt{clobbered} accumulates any session that, while
tearing down, removed state another session had installed.

\texttt{looprun} is the loop's run-to-completion token, as described in
the introduction.

% ============================================================
\section{Initialization}
% ============================================================

\begin{schema}{Init}
  Leg'
\where
  phase' = (\lambda s : SESSION @ sfresh) \\
  listener' = \emptyset \\
  writer' = \emptyset \\
  incumbent' = \emptyset \\
  owner' = \emptyset \\
  gated' = \emptyset \\
  held' = \emptyset \\
  leased' = leaseoff \\
  looprun' = \emptyset \\
  clobbered' = \emptyset
\end{schema}

A single initial state: nothing connected, nothing registered, no lease.

% ============================================================
\section{The Connect Path}
% ============================================================

\texttt{Connect} is \texttt{run}'s prologue, which the loop executes as
one uninterrupted run between the \texttt{accept} await and the
handshake write: the registry entry is recorded (renewing the lease),
the pub-sub writer is bound, and the session installs itself as the
connection's listener, displacing whatever listener was there.

Installing a listener also clears the connection's callbacks.  This is
the rule that makes the listener slot the single owner of the
connection's callback set: a callback may exist only while the session
that registered it holds the slot, so a new occupant starts empty and
the app re-registers from its \texttt{on\_connect}.  Without this
clause a superseded session's entries outlive it with no one left to
withdraw them (Section~\ref{sec:fidelity}, variant~3).

\begin{schema}{Connect}
  \Delta Leg \\
  s? : SESSION
\where
  looprun = \emptyset \\
  phase~s? = sfresh \\
  phase' = phase \oplus \{ s? \mapsto sattached \} \\
  listener' = \{ s? \} \\
  writer' = \{ s? \} \\
  incumbent' = \{ s? \} \\
  owner' = \emptyset \\
  leased' = leaseon \\
  gated' = gated \\
  held' = held \\
  looprun' = looprun \\
  clobbered' = clobbered
\end{schema}

\texttt{Ready} is the handshake write landing and the read loop starting.

\begin{schema}{Ready}
  \Delta Leg \\
  s? : SESSION
\where
  looprun = \emptyset \\
  phase~s? = sattached \\
  phase' = phase \oplus \{ s? \mapsto spumping \} \\
  leased' = leaseon \\
  listener' = listener \\
  writer' = writer \\
  incumbent' = incumbent \\
  owner' = owner \\
  gated' = gated \\
  held' = held \\
  looprun' = looprun \\
  clobbered' = clobbered
\end{schema}

\texttt{Drop} is the loss of the socket --- a peer that went away, a
protocol violation, or a handshake write that failed --- followed by the
suspension this model turns on.  The session is no longer serving
anything, but it has removed nothing: it sits in \texttt{\_pump}'s
\texttt{finally}, awaiting the writer task it just cancelled, with all
its installed state still in place.

\begin{schema}{Drop}
  \Delta Leg \\
  s? : SESSION
\where
  looprun = \emptyset \\
  phase~s? \in \{ sattached, spumping \} \\
  phase' = phase \oplus \{ s? \mapsto ssuspended \} \\
  listener' = listener \\
  writer' = writer \\
  incumbent' = incumbent \\
  owner' = owner \\
  gated' = gated \\
  held' = held \\
  leased' = leased \\
  looprun' = looprun \\
  clobbered' = clobbered
\end{schema}

% ============================================================
\section{The Teardown Path}
% ============================================================

Teardown resumes after the suspension and asks one question first: am I
still the connection's listener?  The answer decides everything, because
the listener slot is the connection's ownership token.  The comparison
and the removals it authorises are one critical section, so no
registration can land between the two.

\texttt{TeardownDetach} is the owned case.  It clears the listener and,
in the same step, the callbacks that listener owned; the ghost
\texttt{clobbered} records a removal of anything installed by another
session, which the guard is there to prevent.  It also takes the loop's
run token, so no reconnect can interleave before the disconnect cascade.

\begin{schema}{TeardownDetach}
  \Delta Leg \\
  s? : SESSION
\where
  looprun = \emptyset \\
  phase~s? = ssuspended \\
  listener = \{ s? \} \\
  phase' = phase \oplus \{ s? \mapsto stearing \} \\
  listener' = \emptyset \\
  owner' = \emptyset \\
  incumbent' = incumbent \setminus \{ s? \} \\
  looprun' = \{ s? \} \\
  clobbered' = \\
  \quad~ \IF listener \cup \ran owner \subseteq \{ s? \} \\
  \quad~ \THEN clobbered \ELSE clobbered \cup \{ s? \} \\
  writer' = writer \\
  gated' = gated \\
  held' = held \\
  leased' = leased
\end{schema}

\texttt{TeardownStale} is the other case: the session is no longer the
listener, so a successor has taken the connection and every piece of
connection-bound state now belongs to that successor.  The stale session
removes nothing and finishes.

\begin{schema}{TeardownStale}
  \Delta Leg \\
  s? : SESSION
\where
  looprun = \emptyset \\
  phase~s? = ssuspended \\
  listener \neq \{ s? \} \\
  phase' = phase \oplus \{ s? \mapsto sdone \} \\
  listener' = listener \\
  writer' = writer \\
  incumbent' = incumbent \\
  owner' = owner \\
  gated' = gated \\
  held' = held \\
  leased' = leased \\
  looprun' = looprun \\
  clobbered' = clobbered
\end{schema}

\texttt{TeardownDisconnect} is the Hub's disconnect cascade --- the
pub-sub writer and the subscriptions --- which runs only after a
successful detach and while the run token is still held, so it cannot
reach a successor's writer.

\begin{schema}{TeardownDisconnect}
  \Delta Leg \\
  s? : SESSION
\where
  phase~s? = stearing \\
  looprun = \{ s? \} \\
  phase' = phase \oplus \{ s? \mapsto sdone \} \\
  writer' = \emptyset \\
  looprun' = \emptyset \\
  clobbered' = \\
  \quad~ \IF writer \subseteq \{ s? \} \THEN clobbered \\
  \quad~ \ELSE clobbered \cup \{ s? \} \\
  listener' = listener \\
  incumbent' = incumbent \\
  owner' = owner \\
  gated' = gated \\
  held' = held \\
  leased' = leased
\end{schema}

A finished session's slot is later occupied by a fresh process.  In the
implementation each reconnect constructs a new \texttt{HubListenSession},
so the reused element must not still be referenced anywhere; the guards
say exactly that.  \texttt{Reset} keeps the system live --- it never runs
out of enabled operations at a benign terminal state --- so that the
deadlock check reports only genuine wedges, never mere termination.

\begin{schema}{Reset}
  \Delta Leg \\
  s? : SESSION
\where
  phase~s? = sdone \\
  looprun = \emptyset \\
  s? \notin listener \\
  s? \notin writer \\
  s? \notin incumbent \\
  s? \notin \ran owner \\
  s? \notin \ran gated \\
  phase' = phase \oplus \{ s? \mapsto sfresh \} \\
  listener' = listener \\
  writer' = writer \\
  incumbent' = incumbent \\
  owner' = owner \\
  gated' = gated \\
  held' = held \\
  leased' = leased \\
  looprun' = looprun \\
  clobbered' = clobbered
\end{schema}

% ============================================================
\section{Registration}
% ============================================================

A registration arrives on another leg entirely --- an MCP tool call or a
REST request from the same identity, resolving to the same connection id
--- and therefore on another thread.  It is two steps because the code
makes it two: the push-reachability gate reads the router, and the
commit writes the client registry.

\texttt{RegisterGate} is \texttt{has\_listener}.  It records which
listener it saw, because that is the fact the commit must still be able
to rely on.

\begin{schema}{RegisterGate}
  \Delta Leg \\
  c? : CALLBACK
\where
  listener \neq \emptyset \\
  leased = leaseon \\
  c? \notin \dom owner \\
  c? \notin \dom gated \\
  gated' = gated \cup (\{ c? \} \cross listener) \\
  phase' = phase \\
  listener' = listener \\
  writer' = writer \\
  incumbent' = incumbent \\
  owner' = owner \\
  held' = held \\
  leased' = leased \\
  looprun' = looprun \\
  clobbered' = clobbered
\end{schema}

\texttt{RegisterCommit} is the registry write.  Its guard is the
compare-and-set that the current code lacks: commit only if the listener
that passed the gate is \emph{still} the listener.  Because the listener
slot and the callback set are under one lock, the comparison and the
write are one critical section, and no teardown can pass between them.

\begin{schema}{RegisterCommit}
  \Delta Leg \\
  c? : CALLBACK
\where
  c? \in \dom gated \\
  listener = \{ gated~c? \} \\
  owner' = owner \oplus \{ c? \mapsto gated~c? \} \\
  gated' = \{ c? \} \ndres gated \\
  phase' = phase \\
  listener' = listener \\
  writer' = writer \\
  incumbent' = incumbent \\
  held' = held \\
  leased' = leased \\
  looprun' = looprun \\
  clobbered' = clobbered
\end{schema}

\texttt{RegisterRefuse} is the same call losing the compare: the leg it
was gated against has gone, so the caller is told it is not
push-reachable rather than given an entry nothing can deliver.

\begin{schema}{RegisterRefuse}
  \Delta Leg \\
  c? : CALLBACK
\where
  c? \in \dom gated \\
  listener \neq \{ gated~c? \} \\
  gated' = \{ c? \} \ndres gated \\
  phase' = phase \\
  listener' = listener \\
  writer' = writer \\
  incumbent' = incumbent \\
  owner' = owner \\
  held' = held \\
  leased' = leased \\
  looprun' = looprun \\
  clobbered' = clobbered
\end{schema}

% ============================================================
\section{Clicks, Lease, and Sweep}
% ============================================================

\texttt{Route} is a click on a registered callback landing in the
connection's hold.  It runs on the click-dispatch thread, so it is not
subject to the run token.

\begin{schema}{Route}
  \Delta Leg \\
  c? : CALLBACK
\where
  c? \in \dom owner \\
  leased = leaseon \\
  held' = held \cup \{ c? \} \\
  phase' = phase \\
  listener' = listener \\
  writer' = writer \\
  incumbent' = incumbent \\
  owner' = owner \\
  gated' = gated \\
  leased' = leased \\
  looprun' = looprun \\
  clobbered' = clobbered
\end{schema}

\texttt{Drain} is the listener taking the hold and pushing its frames.

\begin{schema}{Drain}
  \Delta Leg \\
  s? : SESSION
\where
  looprun = \emptyset \\
  listener = \{ s? \} \\
  phase~s? = spumping \\
  held' = \emptyset \\
  phase' = phase \\
  listener' = listener \\
  writer' = writer \\
  incumbent' = incumbent \\
  owner' = owner \\
  gated' = gated \\
  leased' = leased \\
  looprun' = looprun \\
  clobbered' = clobbered
\end{schema}

Any frame from a pumping session renews the lease.  This is the
operation that makes a broken state permanent rather than self-healing:
a session that is live and keeps renewing keeps the entry alive, so the
sweep that would clear a wrongly-emptied entry never runs.

\begin{schema}{LeaseRenew}
  \Delta Leg \\
  s? : SESSION
\where
  phase~s? = spumping \\
  leased' = leaseon \\
  phase' = phase \\
  listener' = listener \\
  writer' = writer \\
  incumbent' = incumbent \\
  owner' = owner \\
  gated' = gated \\
  held' = held \\
  looprun' = looprun \\
  clobbered' = clobbered
\end{schema}

\begin{schema}{LeaseLapse}
  \Delta Leg \\
\where
  leased = leaseon \\
  \lnot (\exists s : SESSION @ phase~s \in \{ sattached, spumping \}) \\
  leased' = leaseoff \\
  phase' = phase \\
  listener' = listener \\
  writer' = writer \\
  incumbent' = incumbent \\
  owner' = owner \\
  gated' = gated \\
  held' = held \\
  looprun' = looprun \\
  clobbered' = clobbered
\end{schema}

\texttt{Sweep} is the opportunistic reap on a lapsed lease: the entry's
callbacks, its pending registrations, and its hold all go.

\begin{schema}{Sweep}
  \Delta Leg
\where
  leased = leaseoff \\
  owner' = \emptyset \\
  gated' = \emptyset \\
  held' = \emptyset \\
  phase' = phase \\
  listener' = listener \\
  writer' = writer \\
  incumbent' = incumbent \\
  leased' = leased \\
  looprun' = looprun \\
  clobbered' = clobbered
\end{schema}

% ============================================================
\section{Invariants}
\label{sec:inv}
% ============================================================

Four properties must hold across all reachable states.  None is placed
in the \texttt{Leg} schema: a predicate placed there would be enforced
as a guard on every successor, silently disabling any operation that
would break it and thereby masking the very defect this model exists to
catch.

\paragraph{Invariant 1 --- the incumbent stays reachable.}
The session whose connect most recently completed, and which has not
itself detached, holds both the listener slot and the writer slot:
\[
  incumbent \neq \emptyset
  \implies (listener = incumbent \land writer = incumbent).
\]
This is the safety core of ``a live session is never left
push-unreachable''.  It is what the stale-teardown clobber breaks.

\paragraph{Invariant 2 --- callbacks belong to the current listener.}
Every registered callback was installed by the session that holds the
listener slot now:
\[
  \ran owner \subseteq listener.
\]
Since \texttt{\# listener $\leq$ 1}, this says both that no callback is
registered without a listener and that none outlives the leg it was
registered against.  It is what the registration gap breaks.

\paragraph{Invariant 3 --- no cross-session removal.}
No session ever removes state another session installed:
\[
  clobbered = \emptyset.
\]
The ghost is written by exactly the two teardown steps that remove
shared state, so its emptiness is a property of the transition relation
expressed as a property of a state, and ProB reports the offending trace
rather than merely the offending state.

\paragraph{Invariant 4 --- deadlock freedom.}
No reachable state leaves the system unable to progress.  The only
mutual-exclusion device is \texttt{looprun}, taken by
\texttt{TeardownDetach} and released by \texttt{TeardownDisconnect};
since \texttt{TeardownDisconnect}'s guard is implied by the state
\texttt{TeardownDetach} establishes, the token cannot be held forever
and no cyclic wait can form.  ProB discharges this over the whole
reachable state space.

% ============================================================
\section{Verification}
\label{sec:verify}
% ============================================================

Type-checking is a \texttt{fuzz} pass:
\begin{verbatim}
  fuzz -t docs/listen_lifecycle.tex
\end{verbatim}

Invariants 1--3 are state properties, checked by asking ProB whether
their negation is \emph{reachable}.  A goal that is not found means the
invariant holds on every reachable state:
\begin{verbatim}
  # Invariant 1 (must report: goal NOT found)
  probcli docs/listen_lifecycle.tex -model_check -nodead \
    -goal "incumbent /= {} & (listener /= incumbent or writer /= incumbent)" \
    -p DEFAULT_SETSIZE 2

  # Invariant 2 (must report: goal NOT found)
  probcli docs/listen_lifecycle.tex -model_check -nodead \
    -goal "not(ran(owner) <: listener)" \
    -p DEFAULT_SETSIZE 2

  # Invariant 3 (must report: goal NOT found)
  probcli docs/listen_lifecycle.tex -model_check -nodead \
    -goal "clobbered /= {}" \
    -p DEFAULT_SETSIZE 2
\end{verbatim}

Invariant 4 is the deadlock check over the full reachable state space:
\begin{verbatim}
  # Invariant 4 (must report: no deadlock, no invariant violation)
  probcli docs/listen_lifecycle.tex -model_check \
    -p DEFAULT_SETSIZE 2
\end{verbatim}

The last run also discharges the structural constraints of the
\texttt{Leg} schema --- in particular that \texttt{owner} and
\texttt{gated} remain functional and that the occupancy slots stay
singletons --- since ProB carries a state schema's predicate as a
machine invariant.

\paragraph{Results.}
At \texttt{DEFAULT\_SETSIZE 2} the reachable space is 3\,136 states over
14\,545 transitions, every operation covered, no deadlock, and all three
goals reported \emph{not found}.  At \texttt{DEFAULT\_SETSIZE 3} --- three
sessions and three callbacks --- it is 277\,136 states over 1\,885\,113
transitions, again exhaustive, again with no counter-example to any goal.
Two sessions are enough to exhibit every defect of
Section~\ref{sec:fidelity}; three confirm that nothing depends on the
carrier being that small.

% ============================================================
\section{What the Model Asks of the Implementation}
\label{sec:consequences}
% ============================================================

The verified design reduces to one idea: \emph{the listener slot is the
connection's ownership token, and every write to connection-bound state
is a compare-and-set against it}.  Three operations follow, and each
must be a single critical section.

\begin{enumerate}
  \item \textbf{Attach} (\texttt{Connect}).  Install this session as the
    connection's listener and writer, and clear the connection's
    callbacks in the same step.  A new occupant of the slot starts with
    no entries; the app re-registers from its \texttt{on\_connect}.
  \item \textbf{Register} (\texttt{RegisterGate} $+$
    \texttt{RegisterCommit}).  Commit the callback only if the listener
    observed at the gate is still the listener.  The comparison and the
    write must be one critical section --- a re-read that is not atomic
    with the write reinstates the gap.
  \item \textbf{Teardown} (\texttt{TeardownDetach} /
    \texttt{TeardownStale}).  If this session is still the listener,
    clear the listener and the callbacks together, then run the
    disconnect cascade; otherwise remove nothing at all.
\end{enumerate}

Two structural consequences follow, and neither is optional.

The listener slot and the callback set must live under \emph{one} lock.
They are in separate registries today --- the slot in
\texttt{CallbackRouter}, the callbacks on the session in
\texttt{HubClientRegistry} --- and the two locks are deliberately never
nested, so as things stand no single critical section can span them.
Moving the slot onto the session, and leaving the router with the holds
alone, brings all three operations under the registry's existing lock.
It adds no lock, nests nothing, and leaves the router's live-read-then-lock
discipline untouched.

Ownership is decided by comparing the stored listener to \texttt{self},
and a generation counter adds nothing: a \texttt{HubListenSession} is
constructed once per run of the \texttt{/ws} route and installs itself
once, so its identity is already a monotone stamp.  A counter would only
restate that fact in a second place, where it could disagree.  What
would be unsound is a token that \emph{can} recur --- the connection id,
or a bare ``I installed something'' flag --- because a stale holder of
such a token compares equal to the incumbent.

The writer slot needs no separate guard.  It is removed by the
disconnect cascade, which runs inside the same loop run as the detach
that authorised it, so no reconnect can install a successor's writer in
between; the model encodes this as the \texttt{looprun} token, and
Invariant~1 holds with the writer removal unguarded.  That conclusion
depends on \texttt{\_teardown} containing no \texttt{await}, which it
must therefore continue not to.

% ============================================================
\section{Fidelity: reproducing the known defects}
\label{sec:fidelity}
% ============================================================

A model that cannot reproduce the defects it guards against proves
nothing.  Each variant below is the specification above with one named
clause removed; each makes a different invariant reachable, and each
trace below is ProB's, not a hand-written one.  Sessions are written $A$
and $B$ for \texttt{SESSION1} and \texttt{SESSION2}, callbacks $c$ and
$d$.

\subsection*{Variant 1 --- the stale-teardown clobber}

Delete the ownership guard \texttt{listener = \{s?\}} from
\texttt{TeardownDetach}, and delete the \texttt{TeardownStale} schema
(with no guard there is no stale case to divert).  This is the current
\texttt{\_teardown}, which calls \texttt{remove\_listener},
\texttt{withdraw\_callbacks}, and \texttt{on\_disconnect} on the
connection id unconditionally.

The shortest trace to a live, push-unreachable session:

\begin{enumerate}
  \item \texttt{Connect(A)}: $A$ is the listener, the writer, and the
    incumbent.
  \item \texttt{Drop(A)}: $A$'s socket goes; $A$ suspends in the
    \texttt{finally}, awaiting its cancelled writer.  Nothing is
    removed yet.
  \item \texttt{Connect(B)}, \texttt{Ready(B)}: the client reconnects
    after its backoff.  $B$ records, binds the writer, installs itself
    as the listener, becomes the incumbent, and starts pumping.
  \item \texttt{TeardownDetach(A)}: $A$ resumes and, unguarded, removes
    $B$'s listener.
  \item \texttt{TeardownDisconnect(A)}: $A$ drops $B$'s writer.
\end{enumerate}

The resulting state has \texttt{incumbent = \{B\}} with
\texttt{listener} and \texttt{writer} both empty --- Invariant~1 ---
and \texttt{clobbered = \{A\}} --- Invariant~3 --- while $B$ is
\texttt{spumping} and \texttt{leased = leaseon}, so no sweep will ever
repair it.  ProB reaches the same violation by a second route in which
$A$ never drops at all, only is superseded; extending the goal to
\texttt{owner = $\emptyset$} produces the thirteen-step trace in which
$B$'s registered callbacks are withdrawn by $A$'s teardown as well.

Under the ownership guard, step~4 cannot fire: at that point
\texttt{listener = \{B\}}, so $A$ takes \texttt{TeardownStale} and
removes nothing.

\subsection*{Variant 2 --- the registration gap}

Restore the guard of variant~1 and instead delete the compare-and-set
\texttt{listener = \{gated\~{}c?\}} from \texttt{RegisterCommit} (and,
with it, \texttt{RegisterRefuse}, whose guard is its complement).  This
is the current \texttt{register\_callback}, whose \texttt{has\_listener}
read and registry write are separate critical sections.  The trace:

\begin{enumerate}
  \item \texttt{Connect(A)}.
  \item \texttt{Drop(A)}.
  \item \texttt{RegisterGate(c)}: the gate reads \texttt{has\_listener}
    as true --- $A$ has not torn down yet --- and returns.  The calling
    thread is now between the two locks.
  \item \texttt{TeardownDetach(A)}: the listener goes and the callbacks
    are withdrawn.
  \item \texttt{RegisterCommit(c)}: the write lands anyway.
\end{enumerate}

The resulting state has \texttt{owner = \{c $\mapsto$ A\}} with
\texttt{listener = $\emptyset$} --- Invariant~2 --- and the lease still
live.  The withdrawal that would have removed the entry has already
run, so nothing will: a menu item that swallows clicks, which is
exactly what the gate exists to prevent.

That this trace survives the ownership guard is the model's sharpest
result.  Guarding the teardown does not close the registration gap: with
variant~1's guard restored and only this clause missing, ProB still
reaches the violation in five steps, while reporting Invariant~1
unreachable over all 23\,488 states.  The two fixes are independent and
both are required.

\subsection*{Variant 3 --- the superseded registration}

Restore both guards and instead replace \texttt{owner' = $\emptyset$} in
\texttt{Connect} with \texttt{owner' = owner} --- the current
\texttt{add\_listener}, which replaces the listener and leaves the
connection's callbacks alone.  Two sessions of one identity may be live
at once, which the identity-derived connection id permits by design:

\begin{enumerate}
  \item \texttt{Connect(A)}, \texttt{RegisterGate(c)},
    \texttt{RegisterCommit(c)}.
  \item \texttt{Connect(B)}: $B$ replaces $A$ in the listener slot.
\end{enumerate}

Now \texttt{owner = \{c $\mapsto$ A\}} while \texttt{listener = \{B\}}
--- Invariant~2.  $A$'s menu entry is still in the bar and its clicks
are routed to $B$.  This is not a race: it is reachable with no
interleaving at all, and it is why the fixed design makes installing a
listener clear the connection's callbacks.

\subsection*{Variant 4 --- the split detach}

Restore everything and instead split \texttt{TeardownDetach} into its
two writes --- \texttt{TeardownRemoveListener} then
\texttt{TeardownWithdraw} --- as \texttt{\_teardown} has them today,
keeping the ownership guard on the first and the run token across both.
This asks whether the two calls may stay separate once the other three
clauses are in place.  They may not:

\begin{enumerate}
  \item \texttt{Connect(A)}, \texttt{Connect(B)},
    \texttt{RegisterGate(c)}, \texttt{RegisterCommit(c)}: $B$ holds the
    slot and owns $c$.
  \item \texttt{Drop(B)}.
  \item \texttt{TeardownRemoveListener(B)}: the listener is gone; the
    callbacks are not.
\end{enumerate}

That intermediate state has \texttt{owner = \{c $\mapsto$ B\}} with
\texttt{listener = $\emptyset$} --- Invariant~2 --- and it is observable,
because \texttt{route} and \texttt{register} run on other threads and
may read the registry inside the window.  The window is transient, so
this is a narrower fault than variant~2's permanent orphan; but it costs
nothing to close, because the recommended fix already places the
listener slot and the callback set under one lock, where clearing both
is one operation.

Each of the four variants is therefore necessary, and together they are
sufficient: with all four clauses in place the model is exhaustively
clean.

\end{document}
